\chapter{Pure Substances Are Hard to Keep ``Pure''}

\begin{kaichang}
Walk along a supermarket's drinks aisle and you can find at least three kinds of ``water'': mineral water, with calcium, magnesium, and potassium contents printed on the bottle; purified drinking water, whose packaging promises ``multiple stages of purification''; and distilled water, also sold in pharmacies and hardware stores, mostly for humidifiers, car batteries, or laboratories.

Before reading on, guess: which is the ``purest''? If you buy properly produced distilled water, does it contain ``only water and absolutely nothing else''? Is there a final endpoint to purity that we can reach?

Write down your answers. We will return to them at the end of the chapter. You will find that purity is one of chemistry's most frequently used and most easily misunderstood concepts.
\end{kaichang}

\section{Looking ``Uniform'' Is Not Enough}

Recall the two types of change from Chapter 1. When sugar dissolves, it does not become a new substance; it merely ``hides'' in water. Heating it into caramel does turn it into something else. Products of the first kind share a feature: two or more things are \textbf{mixed together}, without any becoming something else. Chemistry gives this distinction names:

\begin{itemize}[itemsep=2pt]
\item \textbf{Pure substance}: something made of only one substance. Distilled water approaches this ideal, as do oxygen, copper wire, and refined salt.
\item \textbf{Mixture}: two or more substances mixed together, with each component retaining its own identity. Air, seawater, milk, and the stainless-steel spoon in your hand are all mixtures.
\end{itemize}

It sounds simple, but deciding whether something is pure or mixed is not. Eyes make unreliable judges. Sugar water is clear and transparent and looks like ``just one thing,'' yet evaporating it leaves sugar at the bottom: two households were there all along. Milk also looks uniform, but it separates into layers after standing for days. Shine light through it and it looks white and cloudy (tiny particles inside scatter light everywhere), quite unlike clear water.

The detective must therefore judge by \textbf{behavior} rather than \textbf{appearance}. Does it separate on standing? Can filter paper trap it? Does evaporation leave a residue? Does it boil at a fixed temperature? These observable, measurable clues establish whether something is pure---just as Chapter 2 taught us to infer invisible things from indirect evidence.

Mixtures themselves fall into two classes. Saltwater, air, and uniform tea, whose composition is the same throughout, are \textbf{homogeneous mixtures} (also called solutions, which are not restricted to liquids). Muddy water, vinaigrette, and granite, with visibly distinct regions, are \textbf{heterogeneous mixtures}. But remember: homogeneous does not mean pure. However uniform saltwater is, water and salt remain two households.

We can now classify our four opening ``suspects'':

\begin{itemize}[itemsep=2pt]
\item \textbf{Air}: a homogeneous gas mixture whose main residents are nitrogen and oxygen, with argon, carbon dioxide, and water vapor too.
\item \textbf{Milk}: a heterogeneous mixture, a jumble of water, fat globules, proteins, lactose, and more, occupying an intermediate region of uniformity.
\item \textbf{Seawater}: a homogeneous mixture, a bowl of soup containing dozens of dissolved salts, chiefly sodium chloride.
\item \textbf{Alloys} (such as stainless steel and brass): homogeneous mixtures of several metals, sometimes with nonmetals, melted together so uniformly that no seams are visible.
\end{itemize}

Another practical observational criterion is that \textbf{pure substances have fixed melting and boiling points; mixtures do not}. Pure ice starts melting at exactly $0\,^{\circ}\mathrm{C}$, and its temperature remains constant throughout melting. Butter, a mixture of fats, gradually softens over a temperature range. Pure water boils at $100\,^{\circ}\mathrm{C}$ (at standard atmospheric pressure), but sugar water's boiling point rises as it boils because escaping water leaves an increasingly concentrated solution. Salting icy roads in winter exploits the tendency of impurities to lower melting points. Chapter 4 will shortly explain the particles behind changes of state.

\section{Zooming In to the Particles}

Macroscopic classification describes appearances. Put on the particle glasses we have used since Chapter 2, and the definitions become sharper:

\textbf{A pure substance contains only one kind of particle.} An ideal glass of pure water contains water molecules throughout, always the same kind; a copper wire contains copper atoms throughout. \textbf{A mixture contains two or more different kinds of particles together, each retaining its identity}: nitrogen molecules, oxygen molecules, argon atoms, and carbon dioxide molecules jostle side by side in air; in seawater, charged particles (ions) from various salts wander among water molecules.

Pure substances can be divided further according to \textbf{how many kinds of atom} their particles contain:

\begin{itemize}[itemsep=2pt]
\item \textbf{Elemental substance}: its particles consist of only one kind of atom. An oxygen molecule joins two oxygen atoms; copper wire contains arrays of copper atoms; diamond contains carbon atoms alone.
\item \textbf{Compound}: different kinds of atoms are firmly joined into a whole in fixed proportions. A water molecule contains two hydrogen atoms and one oxygen atom; a salt crystal is a lattice of alternating sodium and chlorine in a one-to-one ratio (Chapter 18 examines this lattice). Water's hydrogen and oxygen cannot be separated without chemical change. Chapter 1's boundary reappears: separating a mixture requires ``sorting,'' but splitting a compound requires ``surgery.''
\end{itemize}

Restate our suspects in particle language. In milk, enormous fat droplets and long protein chains float in a sea of water molecules. These particles are thousands or tens of thousands of times bigger than molecules, which is why milk scatters light and separates into layers. It occupies the gray area between homogeneous and heterogeneous, called a colloid in chemistry. In stainless steel, iron, chromium, and nickel atoms sit together in the same crystal array, mixed uniformly down to the atomic scale. These examples suggest that uniformity fundamentally depends on the size of the particles being mixed. From a few molecules to droplets visible to the eye, there is a continuous spectrum, not two separate drawers.

One easily overlooked statement deserves emphasis: \textbf{mixing does not change the particles themselves}. A nitrogen molecule among oxygen molecules is still the same nitrogen molecule; a sodium ion in dissolved salt is still the same sodium ion. Mixing means ``living together,'' not ``getting married.'' That is why mixtures have no fixed new properties: saltwater is salty, water is not, and each keeps its own accounts.

\section{Separation: Using Differences, Not Performing Magic}

Since these households are ``living together without getting married,'' separating them requires no breaking. We need only identify \textbf{a difference in some property}, then design an apparatus that magnifies it into ``you go this way, and you go that way.'' This is the chapter's most important rule: \textbf{every separation method exploits a difference in properties}.

\begin{itemize}[itemsep=2pt]
\item \textbf{Filtration}: exploits differences in \textbf{particle size}. Filter paper's pores let water molecules and small dissolved particles through while trapping larger particles. Coffee filters trapping grounds, soy-milk makers straining pulp, and masks stopping droplets use the same trick.
\item \textbf{Distillation}: exploits differences in \textbf{boiling point}. Heating turns lower-boiling components into gas first, which is then cooled and recovered as liquid. Seawater desalination, alcohol distillation, and laboratory distilled-water production all rely on it.
\item \textbf{Extraction}: exploits differences in \textbf{which solvent a substance prefers}. The same solute may dissolve very differently in water and oil. Soaking petals in alcohol to extract fragrance and pigments, or extracting medicinal components from plants with suitable solvents (a similar approach was used for artemisinin), makes target molecules ``move house'' to the side that welcomes them more readily.
\item \textbf{Chromatography}: exploits small differences in \textbf{reluctance to leave home}. A flowing liquid or gas carries a mixture past a stationary surface. Components that like sticking to the surface move slowly; those that prefer the fluid move quickly. The households spread out along a strip. Separating ink colors is chromatography.
\item \textbf{Crystallization}: exploits differences in \textbf{how solubility changes with temperature or water quantity}. In salt ponds, evaporating seawater leaves salt with nowhere to stay, so it lines up into crystals. Hot concentrated sugar water forming sugar crystals as it cools follows the same principle.
\item \textbf{Centrifugation}: exploits differences in \textbf{density and settling speed}. A machine spins rapidly, driving heavier components outward to settle at the bottom. Hospitals use this to separate a blood sample into plasma and blood cells.
\end{itemize}

Let us organize these into a detective's tool table:

\begin{table}[htbp]
\centering
\begin{tabular}{lll}
\toprule
Method & Difference used & Everyday or industrial example \\
\midrule
Filtration & Particle size & Coffee filters, water filters \\
Distillation & Boiling point & Desalination, distilled water \\
Extraction & Solubility in different solvents & Fragrances, drug ingredients \\
Chromatography & Adsorption on a fixed surface & Ink separation, drug testing \\
Crystallization & Solubility vs. temperature, water & Salt ponds, rock sugar \\
Centrifugation & Density and settling speed & Blood separation, spin dryers \\
\bottomrule
\end{tabular}
\end{table}

One final rule is nonnegotiable: \textbf{separation neither creates nor destroys matter}. Separating saltwater into salt and pure water adds or removes no salt particle and loses no water molecule. All the atoms have merely been ``assigned to different rooms.'' Remember this conservation-based accounting idea. In Chapter 26, it becomes the absolute foundation of chemical-equation bookkeeping.

\begin{qiaoliang}
How much is ``just a trace of impurity''? Let us calculate. A bottle of purified drinking water contains about \ud{500}{g}. Qualifying purified water requires impurities (total dissolved solids) below about $10$ ppm. Here ppm means ``parts per million'': at most 10 parts of impurities in a million parts by mass. Even at $10$ ppm, the bottle's impurity mass is only:
\[500\ \mathrm{g} \times \frac{10}{10^{6}} = 5\times 10^{-3}\ \mathrm{g}\]
Five milligrams, invisible to the naked eye. Does that sound utterly pure? Calculate again at the particle scale. A water molecule has a mass of about $3\times 10^{-23}$ g, so $500$ g of water contains about
\[\frac{500}{3\times 10^{-23}} \approx 1.7\times 10^{25}\ \text{water molecules}\]
Impurity particles differ from water molecules by only a few orders of magnitude in mass, so those five milligrams contain about $10^{20}$ particles---one hundred times ten quadrillion. The conclusion is counterintuitive: an amount ``almost zero'' macroscopically is still an enormous crowd at the particle scale. Astronomical numbers separate ``a trace'' from ``none,'' which is why detecting trace pollutants is so difficult and so important. Chapter 5 will provide formal tools for this habit of estimating orders of magnitude before drawing conclusions.
\end{qiaoliang}

\begin{shiyan}
\textbf{Try It: ``Taking Apart'' Black Ink with Paper Chromatography}

Materials: a black or dark water-based pen, filter paper (or absorbent kitchen paper towel), a glass, clean water, a pencil, scissors, and a paper clip or toothpick.

Procedure:
\begin{enumerate}[itemsep=1pt]
\item Cut a paper strip about 2 cm wide and 10 cm long. Lightly draw a pencil line across it about 2 cm from one end (water will not carry the pencil mark along).
\item Make a small, concentrated ink dot at the center of the line and let it dry. For a clearer result, add another dot after the first dries.
\item Pour water into the glass to a depth of about 1 cm. Support the unmarked end of the strip over the glass with a toothpick so that the marked end dips into the water. \textbf{The water level must remain below the pencil line and ink dot.}
\item Leave undisturbed and observe for 10--20 minutes. Water climbs the paper and carries pigments along as it passes through the dot.
\end{enumerate}

What to observe: The ink dot stretches into several differently colored bands, often blue, purple, yellow, and others for black ink, and the colors climb to different heights. Faster-moving colors prefer traveling with water; slower ones prefer the paper surface. Several dyes were living in one pen, though your eyes never told you.

Safety: Using only water and an ordinary water-based pen, this experiment needs no hazardous reagents. Do not submerge the paper's ink dot directly in the water (it will dissolve and spoil the experiment). Wash your hands afterward. If a teacher demonstrates an advanced version with alcohol or other solvents, watch in the laboratory; do not try it at home.
\end{shiyan}

\section{Now Let the Symbols Enter}

With the concepts established, chemistry's shorthand can enter. The rule at the symbolic level is: \textbf{only pure substances have their own chemical formulas}, because a formula describes ``what that one kind of particle is like.''

A water molecule has one oxygen atom holding two hydrogen atoms, written \chem{H_2O}; an oxygen molecule has two oxygen atoms, written \chem{O_2}; carbon dioxide has one carbon atom and two oxygen atoms, written \chem{CO_2}. These are all compounds---their particles contain more than one kind of atom. Oxygen, copper (\chem{Cu}), and diamond (\chem{C}), by contrast, are elemental substances---their particles contain only one kind of atom.

Salt is written \chem{NaCl}, but be careful: it does not contain individual ``sodium chloride molecules.'' It is a crystal lattice of alternating sodium and chloride ions in a one-to-one ratio; \chem{NaCl} records that fixed ratio (Chapter 18 explains it in detail).

What about mixtures? \textbf{No single chemical formula can describe them.} Air has no formula; you can only list its components by volume: nitrogen \chem{N_2} about $78\%$, oxygen \chem{O_2} about $21\%$, argon \chem{Ar} about $0.9\%$, carbon dioxide \chem{CO_2} about $0.04\%$, plus varying amounts of water vapor. Notice what the list reveals: even ``pure air'' mixes four or five gases. Is anything truly pure? That question is the protagonist of our next story.

\section{A New Element Hidden in the Third Decimal Place}

\begin{zhengju}
In the late nineteenth century, chemists thought air was fully understood: remove oxygen and water vapor, and ``pure nitrogen'' remained. In 1892, the British physicist Lord Rayleigh was doing something apparently dull---measuring gas densities precisely---when he encountered trouble. The ``nitrogen'' obtained by removing oxygen, carbon dioxide, and water vapor from air had a density of \ud{1.2572}{g/L}; nitrogen produced chemically by decomposing ammonia had a density of \ud{1.2508}{g/L}. The difference was only $0.5\%$, showing up in the third decimal place.

A careless person might have waved it away as ``experimental error.'' Rayleigh refused. He excluded mundane explanations one by one. Poor calibration? He changed apparatus and repeated measurements; the difference persisted. Residual oxygen or water in atmospheric nitrogen? Experiments removed each, but the difference remained. Lighter hydrogen contaminating chemically produced nitrogen? Several chemical sources all yielded lighter gas. Once he had exhausted his excuses, he did something quintessentially scientific: he published the discrepancy in Nature and invited explanations. He did not hide the ``embarrassing error.''

The chemist William Ramsay took up the puzzle. His alternative explanation sounded absurd at the time: \textbf{the ``nitrogen'' from air was not pure at all}. It contained an unknown gas heavier than nitrogen that reacted with nothing. The test was straightforward: pass atmospheric nitrogen repeatedly over glowing magnesium, which ``consumes'' nitrogen to form magnesium nitride. If an unresponsive component existed, it should remain. Indeed, about $1/80$ of the original gas volume stubbornly refused to react. Ramsay examined its spectrum: its colored lines matched no known element. In 1894, the two announced the new element argon, named from the Greek for ``idle,'' because it refused to react with anything.

Remember three lessons. First, the consensus of ``pure nitrogen'' had been wrong for more than a century, by only about one percent: \textbf{your purity claim can never exceed your measurement precision}. Second, anomalous data are clues, not rubbish; had Rayleigh smoothed away the ``error,'' argon might have waited many more years. Third, the scientific community corrects itself through reproducible numbers, not authority. Anyone measuring accurately enough could see that $0.5\%$. The reason for argon's ``idle'' character awaits the elemental neighborhoods of Chapter 10.
\end{zhengju}

\section{The Limits of ``Pure'': A Continuous Spectrum}

We can now define the boundaries of purity. This chapter's model divides matter into two drawers, pure substances and mixtures. It works in most cases, but has three limits you must understand.

\textbf{Limit one: purity is a continuous scale, not a black-or-white label.} Reagent grades show this: analytical grade is roughly $99.5\%$ or higher, guaranteed reagent grade is stricter, and semiconductor silicon reaches $99.999999999\%$---eleven nines, meaning at most one impurity atom per hundred billion silicon atoms. Why such extremes? A transistor is only a few hundred atoms wide, and impurity atoms deliberately or accidentally alter the paths of current (we will meet this in Chapter 16 and later materials chapters). Purity is thus an engineering question, not a philosophical one: \textbf{whether it is sufficient depends on its use}.

\textbf{Limit two: even ``the same molecule'' hides another layer---isotopes.} All molecules in pure water are indeed water, but hydrogen has two faces. Most is ordinary hydrogen, while about 1 in every 3200 hydrogen atoms is ``heavy hydrogen,'' deuterium, containing a neutron. Pure water therefore naturally contains tiny amounts of ``heavy-water molecules.'' Does this count as impurity? Chemistry usually says no, since isotopes have almost identical chemical personalities (Chapter 8 explains their differences and uses). Yet even ``one hundred percent water'' does not withstand microscopic questioning.

\textbf{Limit three: no knife-edge separates homogeneous from heterogeneous.} Colloids such as milk, soy milk, fog, and blood contain particles much larger than molecules but much smaller than sand, and can stand a long time without separating. Are they homogeneous or heterogeneous? The answer depends on your measuring ruler: uniform to the naked eye, visibly mixed under a microscope. Two drawers cannot hold a continuous world. All classification models share this weakness: useful, but not laws of the world itself.

\begin{wuqu}
``The purer, the better and safer,'' and ``natural extracts are always purer and healthier than synthetic chemicals'': two widespread misconceptions.

Consider counterexamples. Pure oxygen sounds ``clean and premium,'' but makes seemingly mild materials burn fiercely; hospitals control oxygen cylinders more strictly than ordinary medicines. Aconitine and nicotine are both natural plant extracts and can be very pure, yet are extremely toxic. Conversely, synthetic ascorbic acid (vitamin C) and that extracted from lemons are the same molecule; your body cannot distinguish their origins.

Correct the underlying reasoning: \textbf{safety depends on the substance and dose, not on labels such as ``pure'' or ``natural''}. ``All natural'' does not mean pure: essential oils may contain hundreds of molecular species, far more mixed than an analytical-grade reagent. Nor does ``chemical ingredient'' mean harmful: water, salt, and sugar are all chemicals. As for ``distilled water is purest, so it is best for you,'' food can supply the small quantities of minerals found in everyday drinking water. Drinking purified water is fine, but the inference that greater purity means greater health has no scientific basis.
\end{wuqu}

\section{Back to the Three Bottles on the Shelf}

We can now keep our opening promise.

\textbf{Mineral water} is a homogeneous mixture: calcium, magnesium, potassium, and other mineral ions are dissolved among water molecules. The label lists its components; its selling point is precisely its ``impurity.'' \textbf{Purified drinking water} has lost most ions and microorganisms but still contains dissolved atmospheric oxygen and carbon dioxide, and perhaps traces from production equipment. Testing standards define its purity, not an absolute absence of everything but water. \textbf{Distilled water} comes closest to the ideal pure substance: distillation leaves nonvolatile salts in the boiler. Yet it is not the endpoint. Air dissolves after opening, containers release trace substances, and isotopes remain: however much you distill it, a ``heavy-water molecule'' still slips in among every three thousand or so water molecules.

The opening question's answer is therefore that distilled water is purest, but ``only water and nothing else'' does not exist. Purity is a graduated ruler, not a medal you can finally win. Science asks how pure, measured by what method, and whether that is enough for the intended use. In terms of the book's five recurring questions, we have only begun to answer one: what is it made of? How its components connect and arrange themselves occupies the next twenty-odd chapters.

\section{Purification Is Big Business, and Particle Motion Runs It}

Purity is no mere laboratory obsession. The phone in your hand is a triumph of purification industry. Chip silicon starts as quartz sand, is reduced to crude silicon about $98\%$ pure, then undergoes chemical conversions and distillation (distillation again!) to reach eleven nines. Finally, ``zone refining'' slowly sweeps a molten region along a silicon rod, sweeping residual impurities to one end to be cut off. In medicine, chromatography is standard detective equipment: drug manufacturers check each batch's purity and impurity profile, and anti-doping laboratories retrieve banned molecules from athletes' urine at billionths of a gram. Remember that a trace is still a vast crowd at the particle scale, enough to establish a case. Environmental monitoring stations likewise separate and then quantify substances in tap water.

Why do these methods work? Distillation depends on boiling points, determined by particles' ability to escape one another. Crystallization depends on solubility changing with temperature, which fundamentally reflects how vigorously particles move. Centrifugation depends on settling, a race involving density and motion. Chromatography depends on the probability of ``moving on or staying.'' Every clue points to the same hidden boss: \textbf{particles' ceaseless motion}. Next, in Chapter 4, we confront it directly: why can the same substance be solid, liquid, or gas? What are particles doing when its state changes? The workings of dissolution itself will receive their own investigation in Chapter 25.

\begin{tiaozhan}
How can chromatography separate two molecules with nearly identical properties? The secret is repeated amplification. Imagine the paper as $n$ steps. At each step, a molecule either advances with the water or pauses, adsorbed on the paper. Suppose molecule A spends $0.50$ of its time at each step advancing, and B spends $0.45$. At one step, a difference of only $10\%$ cannot separate them. After $n$ steps, however, their average positions are proportional to $0.50n$ and $0.45n$, so the gap grows linearly with $n$. Each band's width grows only as $\sqrt{n}$, following the law of random fluctuations. Separation thus grows as $n$, while blurring grows only as $\sqrt{n}$. With enough steps, even two very similar molecules form cleanly separated bands. Industrial chromatography columns can provide tens of thousands of ``steps,'' the source of their formidable resolving power. The same logic---tiny differences multiplied by enormous repetition---will return in Chapter 29 on reaction rates and later in evolution. Skipping this box does not interrupt the main story.
\end{tiaozhan}

\section*{Try It Yourself}

\begin{sikaoti}
\item Classify the following as pure substances or mixtures, subdividing pure substances into elemental substances or compounds, and draw a particle diagram for each: clean air, sucrose crystals, stainless steel, milk, and oxygen. Add a note that circle sizes do not represent atoms' true relative sizes.
\item A glass of sugar water and a glass of soda are both homogeneous mixtures. Draw particle diagrams identifying at least three kinds of particles in each, then explain in one sentence why homogeneous does not mean pure.
\item You receive a handful of salt mixed with sand and iron filings and want to recover the salt. Design a separation procedure as a material-flow diagram, stating the property difference used at each step. Explain what happens to the total number of atoms throughout.
\item In paper chromatography, a blue band climbs higher than a yellow one. Which pigment prefers moving water, and which prefers the paper? Predict how the gap between bands would change with paper twice as long and explain why (you may use the challenge box's reasoning).
\item If Rayleigh's instruments could measure only to one decimal place, would argon still have been discovered? Use this example to explain the relationship between measurement precision and purity claims.
\item An advertisement describes a product as ``$100\%$ pure, absolutely free of chemicals.'' Refute it from two angles: why purity cannot reach one hundred percent (use isotopes or dissolved gases as an example), and what is wrong with ``absolutely free of chemicals'' itself.
\end{sikaoti}
